%! Periodic Phenomena — Unit 3, Lesson 3.1 — Homework
\documentclass[10pt]{article}
\usepackage{precalculus-article}
\usepackage{precalculus-boxes}
\newcommand{\TallMath}[1]{$\displaystyle #1\rule[-1.4em]{0pt}{3.2em}$}

\begin{document}

\pageheader{Unit 3, Lesson 3.1}{Homework}
\namedateperiod

\vspace{0.10in}

\begin{notesbox}{Problems 1--6}

\textbf{1.}\quad Each table below shows values of a function. Determine whether
each is periodic. If it is, state the period. If it is not, explain why.

\vspace{4pt}
\textbf{(a)}
\begin{center}
\renewcommand{\arraystretch}{1.35}
\begin{tabular}{|c|c|c|c|c|c|c|c|c|}
\hline
$x$ & 0 & 2 & 4 & 6 & 8 & 10 & 12 & 14 \\
\hline
$f(x)$ & 1 & 4 & 7 & 4 & 1 & 4 & 7 & 4 \\
\hline
\end{tabular}
\end{center}
Periodic? \underline{\hspace{1.5cm}} \quad Period $k =$ \underline{\hspace{1.5cm}}
\quad Reason: \writeline

\vspace{8pt}
\textbf{(b)}
\begin{center}
\renewcommand{\arraystretch}{1.35}
\begin{tabular}{|c|c|c|c|c|c|c|c|c|}
\hline
$x$ & 0 & 1 & 2 & 3 & 4 & 5 & 6 & 7 \\
\hline
$g(x)$ & 2 & 5 & 9 & 14 & 20 & 27 & 35 & 44 \\
\hline
\end{tabular}
\end{center}
Periodic? \underline{\hspace{1.5cm}} \quad Reason: \writeline

\vspace{0.12in}
\textbf{2.}\quad A periodic function $h$ has period $k = 5$. You are told that
$h(2) = 13$ and $h(3) = -1$.

\vspace{4pt}
(a)\quad What is $h(7)$? Explain your reasoning.
\writelines{2}

(b)\quad What is $h(12)$? Explain.
\writelines{2}

(c)\quad On what other input values in the domain must $h$ equal $-1$?
\writelines{2}

\vspace{0.12in}
\textbf{3.}\quad \textbf{(Tides Context)} The depth of water at a harbor
entrance (in meters) is approximately periodic. Over one 12-hour period the
depth follows this pattern:

\begin{center}
\renewcommand{\arraystretch}{1.35}
\begin{tabular}{|c|c|c|c|c|c|c|c|}
\hline
Time $t$ (hr) & 0 & 2 & 4 & 6 & 8 & 10 & 12 \\
\hline
Depth $d(t)$ (m) & 3 & 5.5 & 8 & 5.5 & 3 & 0.5 & 3 \\
\hline
\end{tabular}
\end{center}

\vspace{4pt}
(a)\quad What is the period of this tidal pattern? \writeline

(b)\quad What is the maximum depth? When (within one period) does it occur?
\writeline

(c)\quad A ship needs at least 6 m of depth to enter safely. On what
approximate time intervals (within one 12-hour period) is entry safe?
\writelines{2}

(d)\quad If the pattern continues, at what times over a 24-hour day would
the depth equal 8 m?
\writelines{2}

\vspace{0.12in}
\textbf{4. (AP-Style MC)}\quad A function $p$ is periodic with period 6. Which
of the following must be true?

\begin{enumerate}[label=(\Alph*), leftmargin=2em, itemsep=3pt]
  \item $p(6) = 0$
  \item $p(x + 6) = p(x) + 6$ for all $x$
  \item $p(x + 6) = p(x)$ for all $x$ in the domain
  \item The graph of $p$ passes through the origin
  \item $p$ is always increasing on each period
\end{enumerate}

Answer: \underline{\hspace{1.5cm}}

\vspace{0.12in}
\textbf{5. (AP-Style MC)}\quad The table shows values of function $q$. Which
conclusion is best supported by the data?

\begin{center}
\renewcommand{\arraystretch}{1.35}
\begin{tabular}{|c|c|c|c|c|c|c|c|c|c|c|}
\hline
$x$ & 0 & 1 & 2 & 3 & 4 & 5 & 6 & 7 & 8 & 9 \\
\hline
$q(x)$ & 6 & 9 & 6 & 3 & 6 & 9 & 6 & 3 & 6 & 9 \\
\hline
\end{tabular}
\end{center}

\begin{enumerate}[label=(\Alph*), leftmargin=2em, itemsep=3pt]
  \item $q$ is periodic with period 2
  \item $q$ is periodic with period 3
  \item $q$ is periodic with period 4
  \item $q$ is linear
  \item $q$ is not periodic
\end{enumerate}

Answer: \underline{\hspace{1.5cm}}

\end{notesbox}

\vspace{0.08in}

\begin{extensionbox}
\textbf{Extension (Optional):} Function $s$ is periodic with period 10.
You know that $s$ is increasing on the interval $[2, 6]$ with an average
rate of change of $3$ units per input unit on that interval.

\vspace{4pt}
(a)\quad What is the average rate of change of $s$ on the interval $[12, 16]$?
Justify using the definition of a periodic function.

\vspace{4pt}
(b)\quad A student claims: ``Since the average rate of change on $[2,6]$ is 3,
the average rate on $[0, 10]$ must also be positive.'' Is this necessarily
true? Explain.
\writelines{3}
\end{extensionbox}

\vspace{0.08in}

\begin{tcolorbox}[colback=navylight, colframe=navy, arc=2mm,
    left=3mm, right=3mm, top=2mm, bottom=2mm]
\small\textbf{Preview --- Lesson 3.2: Sine, Cosine, and Tangent}\\
We've been identifying and describing periodic relationships from tables.
In Lesson 3.2, we'll meet two functions --- sine and cosine --- that are
the \emph{mathematical model} for the smoothly repeating patterns we saw
today. Come ready to connect period, increase/decrease, and concavity
to the vocabulary of sinusoidal functions.
\end{tcolorbox}

\end{document}
